% Auto-generated by pipeline/scripts/amplicon_length_table.py
% via rule amplicon_length_table
\begin{table}[H]
\centering
\begin{tabular}{lrrrrr}
\toprule
\textbf{Amplicon} & \textbf{Amp. length} & \textbf{Fwd primer} & \textbf{Insert} & \textbf{Rev primer} & \textbf{R1/R2 overlap} \\
\midrule
Amplicon~1 & 195\,bp & 30\,bp & 135\,bp & 30\,bp & 105\,bp \\
Amplicon~2 & 212\,bp & 29\,bp & 153\,bp & 30\,bp & 88\,bp \\
Amplicon~3 & 214\,bp & 30\,bp & 163\,bp & 21\,bp & 86\,bp \\
Amplicon~4 & 237\,bp & 28\,bp & 187\,bp & 22\,bp & 63\,bp \\
\bottomrule
\end{tabular}
\caption{Amplicon composition derived from BLAST primer coordinates (Table~\ref{tab:blast-primers-relaxed}) and primer sequences (Table~\ref{tab:primers}). Insert length $= l_\mathrm{amp} - l_\mathrm{fwd} - l_\mathrm{rev}$. R1/R2 overlap $= 2r - l_\mathrm{amp}$ where $r = 150$\,bp; positive values indicate the two reads overlap in the centre of the amplicon (see Figure~\ref{fig:amplicon-structure}).}
\label{tab:amplicon-lengths}
\end{table}
